\documentclass[11pt]{article}
\usepackage[a4paper,margin=2cm]{geometry}
\usepackage{amsmath,amssymb}
\usepackage[hidelinks]{hyperref}
\allowdisplaybreaks
\title{Qpu.General.Meaning}
\author{Tsvetan Rouschev}
\date{@uuidna/qpu v1.1.1, CC-BY-NC-ND-4.0}
\begin{document}
\maketitle
\noindent\url{https://github.com/uuidna/qpu/blob/v1.1.1/lean/Qpu/General/Meaning.lean}

\paragraph{\texttt{Qpu.gcd\_eq\_nat\_gcd}}
Computes the greatest common divisor of two natural numbers. The GCD of two natural numbers is the largest natural number that evenly divides both.  In particular, the GCD of a number and `0` is the number itself.  This reference implementation via the Euclidean algorithm is overridden in both the kernel and the compiler to efficiently evaluate using arbitrary-precision arithmetic. The definition provided here is the logical model.  Examples: * `Nat.gcd 10 15 = 5` * `Nat.gcd 0 5 = 5` * `Nat.gcd 7 0 = 7`
\[ \forall (a\ b : \mathbb{N}), \mathrm{\_root\_.gcd}\ a\ b = \mathrm{a.gcd}\ b \]
\noindent axioms: \texttt{propext}, \texttt{Quot.sound}; proof: \url{https://github.com/uuidna/qpu/blob/v1.1.1/lean/Qpu/General/Meaning.lean\#L98}

\paragraph{\texttt{chimera\_eq\_mul}}
`a * b` computes the product of `a` and `b`. The meaning of this notation is type-dependent.   Conventions for notations in identifiers:   * The recommended spelling of `*` in identifiers is `mul`.
\[ \forall (m : \mathbb{N}), \mathrm{chimera}\ m = 8 \cdot m ^{2} \]
\noindent axioms: \texttt{propext}; proof: \url{https://github.com/uuidna/qpu/blob/v1.1.1/lean/Qpu/General/Meaning.lean\#L283}

\paragraph{\texttt{choose\_eq\_nat\_choose}}
`choose n k` is the number of `k`-element subsets in an `n`-element set. Also known as binomial coefficients. For the fact that this is the number of `k`-element-subsets of an `n`-element set, see `Finset.card\_powersetCard`.
\[ \forall (n\ k : \mathbb{N}), \mathrm{choose}\ n\ k = \mathrm{n.choose}\ k \]
\noindent axioms: none; proof: \url{https://github.com/uuidna/qpu/blob/v1.1.1/lean/Qpu/General/Meaning.lean\#L49}

\paragraph{\texttt{choose\_three\_succ}}
One more textbook joins every pair already there into a new trinity: `choose (t + 1) 3 = choose t 3 + choose t 2`, for every `t`. Pairs and trinities are one ladder.
\[ \forall (t : \mathbb{N}), \mathrm{choose} (t + 1) 3 = \mathrm{choose}\ t\ 3 + \mathrm{choose}\ t\ 2 \]
\noindent axioms: none; proof: \url{https://github.com/uuidna/qpu/blob/v1.1.1/lean/Qpu/General/Meaning.lean\#L186}

\paragraph{\texttt{choose\_two\_eq\_mixes}}
Mixing textbooks two at a time: `t` textbooks give `t * (t - 1) / 2` mixes, for every `t`. The count the nested catalogue states for its textbooks is this law at their number.
\[ \forall (t : \mathbb{N}), \mathrm{choose}\ t\ 2 = t \cdot (t - 1) / 2 \]
\noindent axioms: \texttt{propext}; proof: \url{https://github.com/uuidna/qpu/blob/v1.1.1/lean/Qpu/General/Meaning.lean\#L150}

\paragraph{\texttt{choose\_two\_succ}}
One more textbook mixes with every textbook already there: `choose (t + 1) 2 = choose t 2 + t`, for every `t`.
\[ \forall (t : \mathbb{N}), \mathrm{choose} (t + 1) 2 = \mathrm{choose}\ t\ 2 + t \]
\noindent axioms: \texttt{propext}; proof: \url{https://github.com/uuidna/qpu/blob/v1.1.1/lean/Qpu/General/Meaning.lean\#L154}

\paragraph{\texttt{cliffordOrder\_eq\_mul}}
`a * b` computes the product of `a` and `b`. The meaning of this notation is type-dependent.   Conventions for notations in identifiers:   * The recommended spelling of `*` in identifiers is `mul`.
\[ \forall (\mathrm{nn} : \mathbb{N}), \mathrm{cliffordOrder}\ \mathrm{nn} = 2 ^{(2 \cdot \mathrm{nn})} \cdot (2 ^{(\mathrm{nn} \cdot \mathrm{nn})} \cdot \prod i \in \mathrm{Finset.range}\ \mathrm{nn}, (2 ^{(2 \cdot (i + 1))} - 1)) \]
\noindent axioms: \texttt{propext}, \texttt{Classical.choice}, \texttt{Quot.sound}; proof: \url{https://github.com/uuidna/qpu/blob/v1.1.1/lean/Qpu/General/Meaning.lean\#L424}

\paragraph{\texttt{factorial\_eq\_nat\_factorial}}
`Nat.factorial n` is the factorial of `n`.
\[ \forall (n : \mathbb{N}), \mathrm{factorial}\ n = \mathrm{n.factorial} \]
\noindent axioms: none; proof: \url{https://github.com/uuidna/qpu/blob/v1.1.1/lean/Qpu/General/Meaning.lean\#L41}

\paragraph{\texttt{fibAux\_eq\_nat\_fib}}
The Fibonacci loop started at consecutive Fibonacci numbers advances the index by its fuel.
\[ \forall (k\ i : \mathbb{N}), \mathrm{fibAux}\ k (\mathrm{Nat.fib}\ i) (\mathrm{Nat.fib} (i + 1)) = \mathrm{Nat.fib} (i + k) \]
\noindent axioms: \texttt{propext}, \texttt{Quot.sound}; proof: \url{https://github.com/uuidna/qpu/blob/v1.1.1/lean/Qpu/General/Meaning.lean\#L122}

\paragraph{\texttt{fibFusion\_eq\_nat\_fib}}
Implementation of the Fibonacci sequence satisfying `fib 0 = 0, fib 1 = 1, fib (n + 2) = fib n + fib (n + 1)`.  *Note:* We use a stream iterator for better performance when compared to the naive recursive implementation.
\[ \forall (n : \mathbb{N}), \mathrm{fibFusion}\ n = \mathrm{Nat.fib}\ n \]
\noindent axioms: \texttt{propext}, \texttt{Quot.sound}; proof: \url{https://github.com/uuidna/qpu/blob/v1.1.1/lean/Qpu/General/Meaning.lean\#L136}

\paragraph{\texttt{fusionStates\_eq\_two\_pow}}
`a \^{} b` computes `a` to the power of `b`. The meaning of this notation is type-dependent.   Conventions for notations in identifiers:   * The recommended spelling of `\^{}` in identifiers is `pow`.
\[ \forall (\mathrm{nn} : \mathbb{N}), \mathrm{fusionStates}\ \mathrm{nn} = 2 ^{\mathrm{nn}} \]
\noindent axioms: \texttt{propext}, \texttt{Quot.sound}; proof: \url{https://github.com/uuidna/qpu/blob/v1.1.1/lean/Qpu/General/Meaning.lean\#L235}

\paragraph{\texttt{gcdAux\_eq\_nat\_gcd}}
Computes the greatest common divisor of two natural numbers. The GCD of two natural numbers is the largest natural number that evenly divides both.  In particular, the GCD of a number and `0` is the number itself.  This reference implementation via the Euclidean algorithm is overridden in both the kernel and the compiler to efficiently evaluate using arbitrary-precision arithmetic. The definition provided here is the logical model.  Examples: * `Nat.gcd 10 15 = 5` * `Nat.gcd 0 5 = 5` * `Nat.gcd 7 0 = 7`
\[ \forall (\mathrm{fuel}\ a\ b : \mathbb{N}), b \le \mathrm{fuel} \to \mathrm{gcdAux}\ \mathrm{fuel}\ a\ b = \mathrm{a.gcd}\ b \]
\noindent axioms: \texttt{propext}, \texttt{Quot.sound}; proof: \url{https://github.com/uuidna/qpu/blob/v1.1.1/lean/Qpu/General/Meaning.lean\#L71}

\paragraph{\texttt{groverIters\_eq\_div}}
`a / b` computes the result of dividing `a` by `b`. The meaning of this notation is type-dependent. * For most types like `Nat`, `Int`, `Rat`, `Real`, `a / 0` is defined to be `0`. * For `Nat`, `a / b` rounds downwards. * For `Int`, `a / b` rounds downwards if `b` is positive or upwards if `b` is negative.   It is implemented as `Int.ediv`, the unique function satisfying   `a \% b + b * (a / b) = a` and `0 ≤ a \% b < natAbs b` for `b ≠ 0`.   Other rounding conventions are available using the functions   `Int.fdiv` (floor rounding) and `Int.tdiv` (truncation rounding). * For `Float`, `a / 0` follows the IEEE 754 semantics for division,   usually resulting in `inf` or `nan`.   Conventions for notations in identifiers:   * The recommended spelling of `/` in identifiers is `div`.
\[ \forall (k : \mathbb{N}), \mathrm{groverIters}\ k = 355 \cdot 2 ^{k} / 452 \]
\noindent axioms: \texttt{propext}, \texttt{Quot.sound}; proof: \url{https://github.com/uuidna/qpu/blob/v1.1.1/lean/Qpu/General/Meaning.lean\#L375}

\paragraph{\texttt{groverSize\_eq\_two\_pow}}
`a \^{} b` computes `a` to the power of `b`. The meaning of this notation is type-dependent.   Conventions for notations in identifiers:   * The recommended spelling of `\^{}` in identifiers is `pow`.
\[ \forall (k : \mathbb{N}), \mathrm{groverSize}\ k = 2 ^{(2 \cdot k)} \]
\noindent axioms: \texttt{propext}, \texttt{Quot.sound}; proof: \url{https://github.com/uuidna/qpu/blob/v1.1.1/lean/Qpu/General/Meaning.lean\#L259}

\paragraph{\texttt{hexbit\_paired\_completely}}
The hexbit paired completely: its 4 qubits form 6 pairs and 4 trinities, each trinity holds 3 pairs, and each pair lies in exactly 2 trinities. The trinities use every pair of the hexbit, each one equally.
\[ \mathrm{hexbit} = 4 \land \mathrm{choose}\ \mathrm{hexbit}\ 2 = 6 \land \mathrm{choose}\ \mathrm{hexbit}\ 3 = 4 \land \mathrm{choose}\ \mathrm{hexbit}\ 3 \cdot \mathrm{choose}\ 3\ 2 = 2 \cdot \mathrm{choose}\ \mathrm{hexbit}\ 2 \]
\noindent axioms: none; proof: \url{https://github.com/uuidna/qpu/blob/v1.1.1/lean/Qpu/General/Meaning.lean\#L220}

\paragraph{\texttt{isingFusion\_eq\_two\_pow}}
`a \^{} b` computes `a` to the power of `b`. The meaning of this notation is type-dependent.   Conventions for notations in identifiers:   * The recommended spelling of `\^{}` in identifiers is `pow`.
\[ \forall (m : \mathbb{N}), \mathrm{isingFusion}\ m = 2 ^{(m - 1)} \]
\noindent axioms: \texttt{propext}, \texttt{Quot.sound}; proof: \url{https://github.com/uuidna/qpu/blob/v1.1.1/lean/Qpu/General/Meaning.lean\#L275}

\paragraph{\texttt{levelGap\_eq\_two\_pow}}
`a \^{} b` computes `a` to the power of `b`. The meaning of this notation is type-dependent.   Conventions for notations in identifiers:   * The recommended spelling of `\^{}` in identifiers is `pow`.
\[ \forall (k : \mathbb{N}), \mathrm{levelGap}\ k = 2 ^{k} \]
\noindent axioms: \texttt{propext}, \texttt{Quot.sound}; proof: \url{https://github.com/uuidna/qpu/blob/v1.1.1/lean/Qpu/General/Meaning.lean\#L243}

\paragraph{\texttt{mint\_eq\_two\_pow}}
`a \^{} b` computes `a` to the power of `b`. The meaning of this notation is type-dependent.   Conventions for notations in identifiers:   * The recommended spelling of `\^{}` in identifiers is `pow`.
\[ \forall (k : \mathbb{N}), \mathrm{mint}\ k = 2 ^{k} \]
\noindent axioms: \texttt{propext}, \texttt{Quot.sound}; proof: \url{https://github.com/uuidna/qpu/blob/v1.1.1/lean/Qpu/General/Meaning.lean\#L35}

\paragraph{\texttt{pair\_in\_trinities}}
Every pair lies in the trinities it forms with each other member: among `t` members a pair is in `t - 2` trinities, so three times the trinities are the pairs times `t - 2`. No pair is left out and none is used more than another.
\[ \forall (t : \mathbb{N}), 3 \cdot \mathrm{t.choose}\ 3 = (t - 2) \cdot \mathrm{t.choose}\ 2 \]
\noindent axioms: \texttt{propext}, \texttt{Quot.sound}; proof: \url{https://github.com/uuidna/qpu/blob/v1.1.1/lean/Qpu/General/Meaning.lean\#L207}

\paragraph{\texttt{paulis\_eq\_two\_pow}}
`a \^{} b` computes `a` to the power of `b`. The meaning of this notation is type-dependent.   Conventions for notations in identifiers:   * The recommended spelling of `\^{}` in identifiers is `pow`.
\[ \forall (\mathrm{nn} : \mathbb{N}), \mathrm{paulis}\ \mathrm{nn} = 2 ^{(2 \cdot \mathrm{nn} + 2)} \]
\noindent axioms: \texttt{propext}, \texttt{Quot.sound}; proof: \url{https://github.com/uuidna/qpu/blob/v1.1.1/lean/Qpu/General/Meaning.lean\#L267}

\paragraph{\texttt{pegasus\_eq\_mul}}
`a * b` computes the product of `a` and `b`. The meaning of this notation is type-dependent.   Conventions for notations in identifiers:   * The recommended spelling of `*` in identifiers is `mul`.
\[ \forall (m : \mathbb{N}), \mathrm{pegasus}\ m = 24 \cdot m \cdot (m - 1) \]
\noindent axioms: none; proof: \url{https://github.com/uuidna/qpu/blob/v1.1.1/lean/Qpu/General/Meaning.lean\#L291}

\paragraph{\texttt{powModAux\_eq}}
The modular-power loop multiplies its accumulator by `a \^{} k`, reduced mod `m`.
\[ \forall (k\ a\ m\ \mathrm{acc} : \mathbb{N}), \mathrm{powModAux}\ k\ a\ m (\mathrm{acc} \bmod m) = \mathrm{acc} \cdot a ^{k} \bmod m \]
\noindent axioms: \texttt{propext}; proof: \url{https://github.com/uuidna/qpu/blob/v1.1.1/lean/Qpu/General/Meaning.lean\#L102}

\paragraph{\texttt{powMod\_eq\_pow\_mod}}
`a \% b` computes the remainder upon dividing `a` by `b`. The meaning of this notation is type-dependent. * For `Nat` and `Int` it satisfies `a \% b + b * (a / b) = a`,   and `a \% 0` is defined to be `a`.   Conventions for notations in identifiers:   * The recommended spelling of `\%` in identifiers is `mod`.
\[ \forall (a\ e\ m : \mathbb{N}), \mathrm{powMod}\ a\ e\ m = a ^{e} \bmod m \]
\noindent axioms: \texttt{propext}; proof: \url{https://github.com/uuidna/qpu/blob/v1.1.1/lean/Qpu/General/Meaning.lean\#L118}

\paragraph{\texttt{qaoaParams\_eq\_mul}}
`a * b` computes the product of `a` and `b`. The meaning of this notation is type-dependent.   Conventions for notations in identifiers:   * The recommended spelling of `*` in identifiers is `mul`.
\[ \forall (p : \mathbb{N}), \mathrm{qaoaParams}\ p = 2 \cdot p \]
\noindent axioms: none; proof: \url{https://github.com/uuidna/qpu/blob/v1.1.1/lean/Qpu/General/Meaning.lean\#L299}

\paragraph{\texttt{qftDepth\_eq\_sub}}
`a - b` computes the difference of `a` and `b`. The meaning of this notation is type-dependent. * For natural numbers, this operator saturates at 0: `a - b = 0` when `a ≤ b`.   Conventions for notations in identifiers:   * The recommended spelling of `-` in identifiers is `sub` (when used as a binary operator).
\[ \forall (q : \mathbb{N}), \mathrm{qftDepth}\ q = 2 \cdot q - 1 \]
\noindent axioms: none; proof: \url{https://github.com/uuidna/qpu/blob/v1.1.1/lean/Qpu/General/Meaning.lean\#L333}

\paragraph{\texttt{qftGates\_eq\_nat\_choose}}
`choose n k` is the number of `k`-element subsets in an `n`-element set. Also known as binomial coefficients. For the fact that this is the number of `k`-element-subsets of an `n`-element set, see `Finset.card\_powersetCard`.
\[ \forall (q : \mathbb{N}), \mathrm{qftGates}\ q = (q + 1)\mathrm{.choose}\ 2 \]
\noindent axioms: \texttt{propext}; proof: \url{https://github.com/uuidna/qpu/blob/v1.1.1/lean/Qpu/General/Meaning.lean\#L145}

\paragraph{\texttt{qftStates\_eq\_two\_pow}}
`a \^{} b` computes `a` to the power of `b`. The meaning of this notation is type-dependent.   Conventions for notations in identifiers:   * The recommended spelling of `\^{}` in identifiers is `pow`.
\[ \forall (q : \mathbb{N}), \mathrm{qftStates}\ q = 2 ^{q} \]
\noindent axioms: \texttt{propext}; proof: \url{https://github.com/uuidna/qpu/blob/v1.1.1/lean/Qpu/General/Meaning.lean\#L251}

\paragraph{\texttt{qftSwaps\_eq\_div}}
`a / b` computes the result of dividing `a` by `b`. The meaning of this notation is type-dependent. * For most types like `Nat`, `Int`, `Rat`, `Real`, `a / 0` is defined to be `0`. * For `Nat`, `a / b` rounds downwards. * For `Int`, `a / b` rounds downwards if `b` is positive or upwards if `b` is negative.   It is implemented as `Int.ediv`, the unique function satisfying   `a \% b + b * (a / b) = a` and `0 ≤ a \% b < natAbs b` for `b ≠ 0`.   Other rounding conventions are available using the functions   `Int.fdiv` (floor rounding) and `Int.tdiv` (truncation rounding). * For `Float`, `a / 0` follows the IEEE 754 semantics for division,   usually resulting in `inf` or `nan`.   Conventions for notations in identifiers:   * The recommended spelling of `/` in identifiers is `div`.
\[ \forall (q : \mathbb{N}), \mathrm{qftSwaps}\ q = q / 2 \]
\noindent axioms: none; proof: \url{https://github.com/uuidna/qpu/blob/v1.1.1/lean/Qpu/General/Meaning.lean\#L324}

\paragraph{\texttt{six\_mul\_choose\_three\_eq\_trinities}}
Mixing textbooks in trinities, three at a time: six times the trinities of `t` textbooks is `t * (t - 1) * (t - 2)`, for every `t`. The count the nested catalogue states for its textbooks is this law at their number.
\[ \forall (t : \mathbb{N}), 6 \cdot \mathrm{choose}\ t\ 3 = t \cdot (t - 1) \cdot (t - 2) \]
\noindent axioms: \texttt{propext}, \texttt{Quot.sound}; proof: \url{https://github.com/uuidna/qpu/blob/v1.1.1/lean/Qpu/General/Meaning.lean\#L181}

\paragraph{\texttt{six\_mul\_nat\_choose\_three}}
Six times the trinities of `t` things is `t * (t - 1) * (t - 2)`, for every `t`: no division.
\[ \forall (t : \mathbb{N}), 6 \cdot \mathrm{t.choose}\ 3 = t \cdot (t - 1) \cdot (t - 2) \]
\noindent axioms: \texttt{propext}, \texttt{Quot.sound}; proof: \url{https://github.com/uuidna/qpu/blob/v1.1.1/lean/Qpu/General/Meaning.lean\#L167}

\paragraph{\texttt{stabProd\_eq\_prod}}
`∏ x ∈ s, f x` is the product of `f x` as `x` ranges over the elements of the finite set `s`.  When the index type is a `Fintype`, the notation `∏ x, f x`, is a shorthand for `∏ x ∈ Finset.univ, f x`.
\[ \forall (k : \mathbb{N}), \mathrm{stabProd}\ k = \prod i \in \mathrm{Finset.range}\ k, (2 ^{(i + 1)} + 1) \]
\noindent axioms: \texttt{propext}, \texttt{Classical.choice}, \texttt{Quot.sound}; proof: \url{https://github.com/uuidna/qpu/blob/v1.1.1/lean/Qpu/General/Meaning.lean\#L384}

\paragraph{\texttt{stabStates\_eq\_mul}}
`a * b` computes the product of `a` and `b`. The meaning of this notation is type-dependent.   Conventions for notations in identifiers:   * The recommended spelling of `*` in identifiers is `mul`.
\[ \forall (\mathrm{nn} : \mathbb{N}), \mathrm{stabStates}\ \mathrm{nn} = 2 ^{\mathrm{nn}} \cdot \prod i \in \mathrm{Finset.range}\ \mathrm{nn}, (2 ^{(i + 1)} + 1) \]
\noindent axioms: \texttt{propext}, \texttt{Classical.choice}, \texttt{Quot.sound}; proof: \url{https://github.com/uuidna/qpu/blob/v1.1.1/lean/Qpu/General/Meaning.lean\#L395}

\paragraph{\texttt{surfaceCorrectable\_eq\_div}}
`a / b` computes the result of dividing `a` by `b`. The meaning of this notation is type-dependent. * For most types like `Nat`, `Int`, `Rat`, `Real`, `a / 0` is defined to be `0`. * For `Nat`, `a / b` rounds downwards. * For `Int`, `a / b` rounds downwards if `b` is positive or upwards if `b` is negative.   It is implemented as `Int.ediv`, the unique function satisfying   `a \% b + b * (a / b) = a` and `0 ≤ a \% b < natAbs b` for `b ≠ 0`.   Other rounding conventions are available using the functions   `Int.fdiv` (floor rounding) and `Int.tdiv` (truncation rounding). * For `Float`, `a / 0` follows the IEEE 754 semantics for division,   usually resulting in `inf` or `nan`.   Conventions for notations in identifiers:   * The recommended spelling of `/` in identifiers is `div`.
\[ \forall (d : \mathbb{N}), \mathrm{surfaceCorrectable}\ d = (d - 1) / 2 \]
\noindent axioms: none; proof: \url{https://github.com/uuidna/qpu/blob/v1.1.1/lean/Qpu/General/Meaning.lean\#L358}

\paragraph{\texttt{surfacePhysical\_eq\_add}}
`a + b` computes the sum of `a` and `b`. The meaning of this notation is type-dependent.   Conventions for notations in identifiers:   * The recommended spelling of `+` in identifiers is `add`.
\[ \forall (d : \mathbb{N}), \mathrm{surfacePhysical}\ d = 2 \cdot d ^{2} - 2 \cdot d + 1 \]
\noindent axioms: \texttt{propext}; proof: \url{https://github.com/uuidna/qpu/blob/v1.1.1/lean/Qpu/General/Meaning.lean\#L341}

\paragraph{\texttt{sympProd\_eq\_prod}}
`∏ x ∈ s, f x` is the product of `f x` as `x` ranges over the elements of the finite set `s`.  When the index type is a `Fintype`, the notation `∏ x, f x`, is a shorthand for `∏ x ∈ Finset.univ, f x`.
\[ \forall (j : \mathbb{N}), \mathrm{sympProd}\ j = \prod i \in \mathrm{Finset.range}\ j, (2 ^{(2 \cdot (i + 1))} - 1) \]
\noindent axioms: \texttt{propext}, \texttt{Classical.choice}, \texttt{Quot.sound}; proof: \url{https://github.com/uuidna/qpu/blob/v1.1.1/lean/Qpu/General/Meaning.lean\#L403}

\paragraph{\texttt{symplectic\_eq\_mul}}
`a * b` computes the product of `a` and `b`. The meaning of this notation is type-dependent.   Conventions for notations in identifiers:   * The recommended spelling of `*` in identifiers is `mul`.
\[ \forall (\mathrm{nn} : \mathbb{N}), \mathrm{symplectic}\ \mathrm{nn} = 2 ^{(\mathrm{nn} \cdot \mathrm{nn})} \cdot \prod i \in \mathrm{Finset.range}\ \mathrm{nn}, (2 ^{(2 \cdot (i + 1))} - 1) \]
\noindent axioms: \texttt{propext}, \texttt{Classical.choice}, \texttt{Quot.sound}; proof: \url{https://github.com/uuidna/qpu/blob/v1.1.1/lean/Qpu/General/Meaning.lean\#L414}

\paragraph{\texttt{thermal\_eq\_mul}}
`a * b` computes the product of `a` and `b`. The meaning of this notation is type-dependent.   Conventions for notations in identifiers:   * The recommended spelling of `*` in identifiers is `mul`.
\[ \forall (\mathrm{millikelvin} : \mathbb{N}), \mathrm{thermal}\ \mathrm{millikelvin} = 1380649 \cdot \mathrm{millikelvin} \cdot 10 \]
\noindent axioms: none; proof: \url{https://github.com/uuidna/qpu/blob/v1.1.1/lean/Qpu/General/Meaning.lean\#L307}

\paragraph{\texttt{trinities\_of\_pairs}}
A trinity is formed by two pairs interacting as one: the two pairs share one member, and through it they are one. Among `t` things, choose the member the pairs share and then the pair of the other two: `t * choose (t - 1) 2` trinities. Every three things form three of them, one for each member the two pairs can share, so the trinities are three times the triples, for every `t`.
\[ \forall (t : \mathbb{N}), t \cdot \mathrm{choose} (t - 1) 2 = 3 \cdot \mathrm{choose}\ t\ 3 \]
\noindent axioms: \texttt{propext}; proof: \url{https://github.com/uuidna/qpu/blob/v1.1.1/lean/Qpu/General/Meaning.lean\#L193}

\paragraph{\texttt{trinity\_pairs\_differ}}
The two pairs of a trinity are two of the pairs there are, and they are different pairs: a trinity on `a`, the shared `b`, and `c` is the pairs `\{a, b\}` and `\{b, c\}`, which differ exactly when `a` and `c` do.
\[ \forall (a\ b\ c : \mathbb{N}), a \ne c \to (a, b) \ne (c, b) \]
\noindent axioms: none; proof: \url{https://github.com/uuidna/qpu/blob/v1.1.1/lean/Qpu/General/Meaning.lean\#L201}

\paragraph{\texttt{two\_mul\_nat\_choose\_two}}
Twice the pairs of `t` things is `t * (t - 1)`, for every `t`: no division.
\[ \forall (t : \mathbb{N}), 2 \cdot \mathrm{t.choose}\ 2 = t \cdot (t - 1) \]
\noindent axioms: \texttt{propext}; proof: \url{https://github.com/uuidna/qpu/blob/v1.1.1/lean/Qpu/General/Meaning.lean\#L158}
\end{document}
